\section{Алгоритмы постобработки и интеграции с CAD-ядром КОМПАС-3D}

После получения сегментации от нейросети (раскрашенный OBJ-файл) выполняется постобработка, подготавливающая данные для передачи в CAD-ядро.

\subsection{Фильтрация треугольников}
Исходная сетка может содержать выбросы – одиночные треугольники, ошибочно отнесённые к операции. Для их удаления применяются следующие методы:
\begin{itemize}
    \item \textbf{Удаление по площади}: удаляются $k\%$ треугольников с наибольшей площадью (параметр \texttt{--remove\_largest\_area}).
    \item \textbf{Кластеризация}: из всех треугольников с меткой 1 оставляется только наибольший связный компонент (флаг \texttt{-c}).
\end{itemize}

\subsection{Преобразование в формат CAD}
Очищенная геометрия сохраняется в файл \texttt{test\_0.stl}. Этот файл затем загружается в КОМПАС-3D, где по нему восстанавливается операция:
\begin{itemize}
    \item Для операций выдавливания (\texttt{bossExtrusion}, \texttt{cutExtrusion}) определяется направление и глубина.
    \item Для скруглений (\texttt{fillet}) и фасок (\texttt{chamfer}) вычисляется радиус или размер фаски.
\end{itemize}

\subsection{Интеграция с КОМПАС-3D}
Разработанный модуль \texttt{main.exe} взаимодействует с КОМПАС-3D через внешний API. Результаты работы (файлы \texttt{meta.json} и \texttt{test\_0.stl}) передаются в среду, где выполняется проверка гипотезы и, при успехе, упрощение модели. Процесс может повторяться итеративно, постепенно восстанавливая дерево построения.

